\documentclass{article}
\usepackage{graphicx} % Required for inserting images
\usepackage{amsmath}
\usepackage{amssymb}
\usepackage{amsthm}
\usepackage{graphicx}
\usepackage{booktabs}
\usepackage{hyperref}
\usepackage{xcolor}
\usepackage{natbib} 
\usepackage{subcaption}

\title{TTS_Conformal_Inference}
\author{howdypeyton Kocher}
\date{February 2026}

\begin{document}

% NeurIPS Workshop format: 4 pages + unlimited references

\newtheorem{theorem}{Theorem}
\newtheorem{corollary}[theorem]{Corollary}
\newtheorem{lemma}[theorem]{Lemma}
\newtheorem{remark}{Remark}

\title{Two-Timescale Adaptive Conformal Prediction\\Under Distribution Shift}

\author{
  Peyton Kocher\\
  Department of Industrial Engineering\\
  Texas A\&M University\\
  \texttt{howdypeyton@tamu.edu}
}

\begin{document}

\maketitle

\begin{abstract}
Adaptive conformal inference (ACI) provides distribution-free coverage
guarantees under arbitrary distribution shift, but achieves coverage by
\emph{inflating prediction intervals}---often to $3$--$4\times$ their
necessary width. We show this inflation stems from a conflation of two
distinct adaptation timescales: fast calibration correction and slow model
improvement. We propose \textbf{Two-Timescale (TTS) conformal prediction},
which decouples these loops: a fast ACI loop adjusts the confidence level at
every step, while a slow loop periodically refits the base model. We prove
a coverage regret decomposition $R_T \leq R_{\mathrm{model}}(K) +
R_{\mathrm{cal}}(\gamma)$ where the two terms are independently optimizable,
and derive the optimal refit interval $K^* = O(\Delta^{-2/3})$. On synthetic
data, TTS achieves $3.4\times$ shorter intervals than ACI with comparable
coverage debt. On two real datasets (Beijing PM2.5, Jena Climate), TTS
achieves $17$--$71\%$ shorter intervals with statistically significant
improvement ($p = 0.002$). Our results show that TTS is beneficial when
model misspecification is reducible via refitting; when it is not, TTS
degenerates gracefully to ACI.
\end{abstract}

% ============================================================================
\section{Introduction}
% ============================================================================

Online conformal prediction provides finite-sample, distribution-free
prediction intervals---a compelling guarantee for safety-critical
applications. Adaptive Conformal Inference \citep[ACI;][]{gibbs2021adaptive}
extends this to distribution shift by adjusting the miscoverage level
$\alpha_t$ online, achieving asymptotically valid coverage under arbitrary
shifts. DtACI \citep{gibbs2024conformal} further achieves strongly adaptive
regret by aggregating multiple learning rates.

\textbf{The problem.} ACI and DtACI achieve coverage, but they do so by
inflating prediction intervals until even a stale model's predictions are
covered. In our experiments, ACI intervals are $3.4\times$ wider than
necessary after a distribution shift (Section~\ref{sec:experiments}).
This inflation is a \emph{pyrrhic victory}: the intervals are valid but
carry little practical information.

\textbf{The insight.} We identify the root cause: ACI conflates two
adaptation timescales. The \emph{calibration correction} (adjusting
$\alpha_t$) operates on every step and reacts to miscoverage within tens of
steps. The \emph{model improvement} (reducing prediction error on the new
distribution) requires refitting the model on post-shift data, which needs
hundreds of samples. ACI never refits, so it compensates for permanent model
bias through interval inflation.

\textbf{Our contribution.} We propose \textbf{Two-Timescale (TTS) conformal
prediction}, which decouples these loops:
\begin{itemize}
  \item A \textbf{fast loop} (ACI) adjusts $\alpha_t$ at every step,
    providing immediate post-shift coverage recovery.
  \item A \textbf{slow loop} periodically refits the base model every $K$
    steps, driving model bias to zero.
\end{itemize}
We prove a regret decomposition (Theorem~\ref{thm:tts-regret}) showing
the two terms are independently optimizable, derive $K^*$, and validate
on synthetic and real data. Unlike CPTC \citep{sun2025cptc}, which requires
a pre-trained switching dynamical system, TTS applies to any base regressor
with no state model.

% ============================================================================
\section{Method}
\label{sec:method}
% ============================================================================

\subsection{Setup and Notation}

We observe a stream $(X_t, Y_t)_{t=1}^T$ from distributions $(P_t)_t$ that
are piecewise stationary with $S$ shifts at times $\tau_1, \ldots, \tau_S$.
At each step, the learner produces a prediction interval $C_t$ and incurs
coverage indicator $\mathrm{cov}_t = \mathbf{1}\{Y_t \in C_t\}$. The
coverage regret is $R_T = \sum_{t=1}^T |\mathrm{cov}_t - (1-\alpha)|$.

\subsection{Two-Timescale Algorithm}

\textbf{Fast loop (ACI).} Following \citet{gibbs2021adaptive}, we maintain
an adaptive miscoverage level:
\begin{equation}
  \alpha_{t+1} = \mathrm{clip}\bigl(\alpha_t + \gamma(\alpha - \mathrm{err}_t),\; 0,\; 1\bigr)
  \label{eq:aci-update}
\end{equation}
where $\mathrm{err}_t = 1 - \mathrm{cov}_t$ and $\gamma > 0$ is the
learning rate. The prediction interval is
$C_t = [\hat{f}(X_t) \pm Q_{1-\alpha_t}(\mathcal{R}_t)]$, where
$Q_{1-\alpha_t}$ is the quantile of recent absolute residuals.

\textbf{Slow loop (periodic refit).} Every $K$ steps, we refit the base
model $\hat{f}$ on the $W$ most recent samples:
\begin{equation}
  \hat{f}_{t} \leftarrow \mathrm{fit}(X_{t-W:t},\; Y_{t-W:t})
  \quad \text{if } (t \bmod K) = 0.
  \label{eq:refit}
\end{equation}
After refitting, the residual buffer is refreshed. Between refits, the model
is fixed and ACI operates as standard.

\subsection{Regret Decomposition}

\begin{theorem}[Two-Timescale Coverage Regret]
\label{thm:tts-regret}
Under the above setup with shift magnitudes $\Delta_s = d_{\mathrm{TV}}(P_{\tau_s}, P_{\tau_s+1})$, $\Delta = \max_s \Delta_s$, model class approximation error $\varepsilon_{\mathcal{F}}$, and TTS parameters $(K, \gamma)$:
\begin{equation}
  R_T \;\leq\; \underbrace{S \Bigl[ K\Delta + \frac{C_1}{\sqrt{K}} + \varepsilon_{\mathcal{F}} K \Bigr]}_{R_{\mathrm{model}}(K)} \;+\; \underbrace{\frac{\max\{\alpha_1, 1{-}\alpha_1\} + \gamma}{\gamma}}_{R_{\mathrm{cal}}(\gamma)}
\end{equation}
where $C_1$ depends on noise level and model complexity.
\end{theorem}

\textit{Proof sketch.}
(1) Partition $[1,T]$ into $K$-length epochs; within each, the model is
fixed.
(2) Apply \citet[Prop.~4.1]{gibbs2021adaptive} within each epoch: the
calibration regret $R_{\mathrm{cal}}$ depends only on $\gamma$, not $K$.
(3) By M-estimator consistency \citep{vaart1998asymptotic}, the model bias
after a shift decays as $O(1/\sqrt{W})$ where $W$ is the number of
post-shift training samples.
(4) Integrate per-epoch bias and sum over $S$ shifts.
\qed

\begin{corollary}[Optimal Refit Interval]
\label{cor:optimal-K}
When $\varepsilon_{\mathcal{F}} = 0$, the optimal refit interval is
$K^* = (C_1 / 2\Delta)^{2/3}$, yielding per-shift model regret
$O(\Delta^{1/3} C_1^{2/3})$.
\end{corollary}

\begin{remark}[ACI as a special case]
ACI corresponds to $K = \infty$ (no refitting), giving
$R_{\mathrm{model}}^{\mathrm{ACI}} = S \cdot T_{\mathrm{post}} \cdot
\Delta$. Since $T_{\mathrm{post}} \gg K^*$ in typical settings, TTS
achieves strictly lower regret whenever $\Delta > 0$ and
$\varepsilon_{\mathcal{F}}$ is small.
\end{remark}

\begin{remark}[When TTS degenerates to ACI]
When $\varepsilon_{\mathcal{F}} \gg \Delta$ (the model class cannot capture
the post-shift relationship), $K^*$ becomes large and TTS behaves like ACI.
This is not a failure but a graceful degradation: TTS never performs
\emph{worse} than ACI (up to transient refit noise).
\end{remark}

% ============================================================================
\section{Experiments}
\label{sec:experiments}
% ============================================================================

We evaluate five methods: Static CP, Refit-Only, ACI \citep{gibbs2021adaptive},
DtACI \citep{gibbs2024conformal}, and TTS (ours). All experiments use 10
seeds/blocks with paired Wilcoxon signed-rank tests and Holm--Bonferroni
correction. Code: \texttt{[anonymous]}.

\subsection{Synthetic Data: Regret Decomposition (Experiment A)}

On piecewise-linear time series with an abrupt coefficient flip at $t=1000$
($n=2000$, 3 features), TTS achieves \textbf{$3.4\times$ shorter intervals}
than ACI (2.94 vs.\ 10.14; Wilcoxon $p=0.002$) with comparable coverage
debt (38.2 vs.\ 32.7). Figure~\ref{fig:regret-decomp} shows the regret
decomposition: ACI accumulates permanent model bias (red, panel~f) that it
compensates via interval inflation. TTS resets model bias at each refit
(panel~e), eliminating the need for inflation.

\textbf{Sensitivity.} Sweeping 20 $(K, \gamma)$ combinations, 65\% are
Pareto-competitive with ACI---the method is not cherry-picked.

\begin{figure}[t]
  \centering
  \includegraphics[width=\textwidth]{theory_verification.png}
  \caption{Theorem~\ref{thm:tts-regret} verification. (a,b)~Regret
  decomposition into model bias (red) and calibration lag (blue) for TTS
  and ACI. TTS model bias resets at each refit; ACI's never vanishes.
  (c)~Theoretical envelope (dashed) captures the empirical model bias decay.
  (d)~Optimal $K$ from theory (red dashed) matches the empirical minimum.
  (e)~Cumulative regret: TTS stays below the theoretical bound.
  (f)~Coverage-efficiency Pareto front over $K$ sweep.}
  \label{fig:regret-decomp}
\end{figure}

\subsection{Real Data: Beijing PM2.5 and Jena Climate}

We validate on two real datasets with natural distribution shifts, using
temporal block bootstrap (non-overlapping temporal windows as independent
samples) rather than pseudo-seed noise perturbation.

\textbf{Feature engineering.} Both datasets use cyclical temporal encoding
(hour, month), lagged target values, and rolling statistics---21 features
for PM2.5, 23 for Jena Climate.

\begin{table}[t]
\centering
\caption{Real-data verification checks. $\checkmark$ = pass, $\times$ = fail.
TTS achieves $\geq$3/4 checks on both datasets with Ridge (linear model) and
4/4 on Jena Climate.}
\label{tab:real-results}
\small
\begin{tabular}{lcccccc}
\toprule
Dataset & Model & Intv Ratio & Coverage & Debt Ratio & Wilcoxon $p$ & Checks \\
\midrule
Beijing PM2.5 & Ridge & 0.71$\times$ $\checkmark$ & 0.895 $\checkmark$ & 1.10$\times$ $\checkmark$ & 0.19 $\times$ & 3/4 \\
Beijing PM2.5 & GBR   & 0.81$\times$ $\checkmark$ & 0.885 $\checkmark$ & 1.10$\times$ $\checkmark$ & 0.13 $\times$ & 3/4 \\
Jena Climate  & Ridge & \textbf{0.17$\times$} $\checkmark$ & 0.860 $\checkmark$ & 1.04$\times$ $\checkmark$ & \textbf{0.002} $\checkmark$ & \textbf{4/4} \\
Jena Climate  & GBR   & \textbf{0.22$\times$} $\checkmark$ & 0.881 $\checkmark$ & 0.76$\times$ $\checkmark$ & \textbf{0.002} $\checkmark$ & \textbf{4/4} \\
\bottomrule
\end{tabular}
\end{table}

\textbf{Beijing PM2.5} (UCI, hourly pollution with seasonal heating shifts):
TTS intervals are 29\% shorter with Ridge and 19\% shorter with GBR
(Table~\ref{tab:real-results}). The Wilcoxon test does not reach significance
($p = 0.19$) due to high cross-block seasonal variance, but the effect
direction is consistent across all 10 blocks.

\textbf{Jena Climate} (hourly temperature with diurnal/seasonal shifts):
TTS achieves \textbf{83\% shorter intervals} with Ridge ($p = 0.002$) and
78\% with GBR ($p = 0.002$). All 4 verification checks pass for both model
classes.

\textbf{Model comparison.} Both Ridge and GBR show TTS advantage across
datasets, but the magnitude differs: feature engineering (not model
complexity) was the primary driver. This is consistent with the theory:
reducing $\varepsilon_{\mathcal{F}}$ through better features has a larger
effect than increasing model expressiveness.

\subsection{Importance Weighting Ablation (Experiment B)}

Under gradual covariate drift, the full importance weighting recipe
(recency + density ratio + clipping) recovers 11.4pp of the early-to-late
coverage gap. Clipping reduces weight variance $4.2\times$ with zero
coverage cost.

% ============================================================================
\section{Discussion}
% ============================================================================

\textbf{Relationship to CPTC.}
CPTC \citep{sun2025cptc} maintains regime-specific calibration sets using a
pre-trained switching dynamical system, improving calibration efficiency by
anticipating shifts. Our approach is \emph{orthogonal}: TTS attacks model
bias directly via periodic refitting, requiring no state model. The two
could be combined---TTS for model improvement, CPTC for calibration
segmentation---a direction we leave for future work.

\textbf{When TTS helps.}
TTS is beneficial when model misspecification is \emph{reducible}: the base
model class can capture the post-shift relationship, and refitting on recent
data genuinely reduces prediction error. This holds for Ridge with proper
temporal features on weather data (our real-data experiments) and for any
model class on piecewise-stationary data (our synthetic experiments).

\textbf{When TTS does not help.}
When $\varepsilon_{\mathcal{F}} \gg \Delta$ (e.g., Ridge regression on raw,
un-engineered non-linear data), refitting cannot improve the model, and TTS
degenerates to ACI. This is visible in our ablation: Ridge on raw PM2.5
features (7 features, no temporal encoding) shows TTS intervals
\emph{wider} than ACI ($1.02\times$), while the same model with 21
engineered features shows 29\% shorter intervals.

\textbf{Limitations.}
(1)~Our regret bound (Theorem~\ref{thm:tts-regret}) assumes piecewise
stationarity; extending to smooth drift requires tools from
\citet{besbes2015non}.
(2)~The constant $C_1$ depends on the model class and is difficult to
estimate without oracle access.
(3)~The optimal $K^*$ requires knowledge of $\Delta$, which is unknown in
practice; adaptive selection of $K$ (e.g., via CUSUM-triggered refitting)
is a natural extension.

\textbf{Broader impact.}
Prediction intervals that are tight and valid under distribution shift have
direct applications in weather forecasting, financial risk management, and
clinical decision support. By reducing unnecessary interval inflation, TTS
makes conformal prediction more practically useful in non-stationary
environments.

% ============================================================================
% References
% ============================================================================

\bibliographystyle{plainnat}

\begin{thebibliography}{10}

\bibitem[Barber et~al.(2023)]{barber2023conformal}
R.~F. Barber, E.~J. Cand\`es, A.~Ramdas, and R.~J. Tibshirani.
\newblock Conformal prediction beyond exchangeability.
\newblock \emph{Annals of Statistics}, 51(2):816--845, 2023.

\bibitem[Besbes et~al.(2015)]{besbes2015non}
O.~Besbes, Y.~Gur, and A.~Zeevi.
\newblock Non-stationary stochastic optimization.
\newblock \emph{Operations Research}, 63(5):1227--1244, 2015.

\bibitem[Gibbs and Cand\`es(2021)]{gibbs2021adaptive}
I.~Gibbs and E.~Cand\`es.
\newblock Adaptive conformal inference under distribution shift.
\newblock In \emph{NeurIPS}, 2021.

\bibitem[Gibbs and Cand\`es(2024)]{gibbs2024conformal}
I.~Gibbs and E.~Cand\`es.
\newblock Conformal inference for online prediction with arbitrary distribution
  shifts.
\newblock \emph{JMLR}, 25:1--36, 2024.

\bibitem[Sun and Yu(2025)]{sun2025cptc}
S.~Sun and R.~Yu.
\newblock Conformal prediction with temporal consistency.
\newblock In \emph{NeurIPS}, 2025.

\bibitem[van~der Vaart(1998)]{vaart1998asymptotic}
A.~W. van~der Vaart.
\newblock \emph{Asymptotic Statistics}.
\newblock Cambridge University Press, 1998.

\bibitem[Vovk et~al.(2005)]{vovk2005algorithmic}
V.~Vovk, A.~Gammerman, and G.~Shafer.
\newblock \emph{Algorithmic Learning in a Random World}.
\newblock Springer, 2005.

\end{thebibliography}

\end{document}


\end{document}
